\chapter{World Model as Data Generator}
\label{chap:mpc-distillation}

\section{Motivation}
\label{sec:mpcd-motivation}

Chapter~\ref{chap:wm-planner} defined planners that use the world model
only at inference time. This keeps model use local and critic-guided, but
it also means that the final policy still depends on test-time search.
This raises a natural question: can the behaviour of the planner be
\emph{absorbed into the policy itself}, so that the final agent
executes at the same level without requiring any online search,
simulation, or test-time optimisation?
Instead of treating the world model as an external decision-time module,
we use it to generate supervision during training, transferring MPC
improvements into the actor so that planning becomes implicit in the
policy parameters.

We explore three progressively more integrated formulations. The
first is a static relabel-and-train pipeline
(\S\ref{sec:mpcd-offline}), where MPC is run once to generate a fixed
dataset of improved actions. The second is an interleaved procedure
(\S\ref{sec:mpcd-interleaved}), where MPC labels and policy updates
are continuously refreshed as training proceeds. The third is a full
online phase inside the world model
(\S\ref{sec:mpcd-online-wm}), where both MPC-generated actions and
world-model-imagined transitions are injected into the replay buffer,
creating generated transition data for continued training.
Across these variants, the central design question is where WM-derived
information enters learning: only as fixed labels, as periodically
updated supervision, or as full synthetic interaction data.

% ---------------------------------------------------------------------
\section{Offline Policy Learning from MPC-Generated Actions}
\label{sec:mpcd-offline}

We begin with the simplest way of converting inference-time gains into
training data: run MPC once over the offline dataset, then train the
policy to imitate the resulting actions.

\paragraph{Phase A --- relabel.}
For every latent state $z \in \mathcal{Z}_{\text{off}}$, we compute a
corresponding improved action chunk $\mathbf{a}^{\text{MPC}}(z)$ by
running Grad-MPC with parameters $N = 32$, $H = 2$, $K = 5$,
$\eta = 0.01$, and Q-term scoring, using the fully trained actor,
critic, and world model.
This produces a fixed relabelled dataset in which each state is paired
with an action that is locally improved according to the planner.

\paragraph{Phase B --- train.}
The BC-flow actor is then trained on these relabelled pairs
$(z, \mathbf{a}^{\text{MPC}})$ using the standard flow-matching loss
from Eq.~\eqref{eq:bc-flow-bg}. During this phase, both the critic
and the original network heads are frozen; only the actor is
updated. Training uses AdamW with learning rate $10^{-4}$, weight
decay $10^{-4}$, batch size $256$, for $10^5$ gradient steps.

\paragraph{Phase C --- evaluate.}
After training, the policy can be evaluated directly, or
used again as the proposal distribution inside MPC (BoN, Grad-MPC, or
FMQ), without requiring any additional training.

\paragraph{Limitations of static action relabelling.}
The main limitation is that the relabelled dataset is generated from a
fixed actor and critic, and therefore reflects a single snapshot of the
optimisation process.
As training progresses, the policy shifts away from the behaviour that
generated the labels, and the frozen relabelling critic cannot improve
the scoring function that defines $\mathbf{a}^{\text{MPC}}$. The
supervision signal therefore becomes stale, limiting continued training.
The next section addresses this by making MPC relabelling part of the
training loop itself.

% ---------------------------------------------------------------------
\section{Interleaved Policy Learning from MPC-Generated Actions}
\label{sec:mpcd-interleaved}

The key limitation of static relabelling is that MPC labels are computed
only once, even though they depend on the actor that proposes actions
and the critic that scores them. A more consistent approach is therefore
to recompute labels periodically as both networks evolve.

\paragraph{Mechanism.}
We maintain a buffer $\mathcal{B}_{\text{MPC}}$ of relabelled
transitions $(z, \mathbf{a}^{\text{MPC}}(z))$. After an initial warm-up
period of $T_{\text{warm}}$ steps, this buffer is refreshed every
$T_{\text{relabel}}$ gradient updates by sampling states from the
offline dataset and recomputing their MPC actions using the current
actor and critic.
During training, each minibatch is formed by mixing real offline
samples with relabelled MPC samples. Specifically, a fraction of each
batch is replaced with entries from $\mathcal{B}_{\text{MPC}}$, and
the standard FQL update is applied to the combined batch.

\paragraph{What the critic sees.}
Importantly, the replacement affects only the state-action pairs. The
rewards and next states remain those of the original offline
transition. Thus, the critic is updated on tuples of the form
$(z, \mathbf{a}^{\text{MPC}}, r, z')$, where $r$ and $z'$ correspond
to the environment transition induced by the original action, not the
MPC action.
In practice, this mismatch is often small because MPC actions remain
close to the behaviour policy due to local refinement. However, it is
conceptually important: the critic is being trained on transitions
that did not actually occur. The fully online variant in
\S\ref{sec:mpcd-online-wm} removes this inconsistency entirely.

\paragraph{Critic update.}
We compare an actor-only variant with a critic-update variant. In the
actor-only setting, MPC labels are used only for an additional policy
loss, while the critic continues to see the original offline actions. In
the critic-update setting, MPC-labelled actions are also inserted into
critic batches. This tests whether the actor can benefit from planner
labels when the critic has never been calibrated on the same action
distribution.

\paragraph{Grad-MPC vs FMQ labels.}
The relabelling procedure can use either Grad-MPC or FMQ. Grad-MPC
produces strong improvements but can introduce noisy action updates due
to iterative optimisation. FMQ produces more stable and structured
updates by constraining refinement to a trust region.
The Grad-MPC and FMQ relabelling results are reported in
\S\ref{sec:exp-mpc-interleaved-grad}--\ref{sec:exp-fmq-distil-inf}
(Tables~\ref{tab:mpc-interleaved-grad}, \ref{tab:mpc-relabel-sweep},
and~\ref{tab:mpc-fmq-distilled}), with the full hyperparameter sweep in
Appendix~\ref{app:oo-sweeps}. The next section removes the remaining
real-transition inconsistency by generating the successor state with the
same world model used for action selection.

\begin{algorithm}[h]
\caption{Interleaved training with MPC-generated actions.}
\label{alg:mpc-interleaved}
\begin{algorithmic}[1]
\Require offline dataset $\mathcal{D}_{\text{off}}$, FQL agent $(\pi_\theta, Q_\phi)$, frozen WM $f_\psi$, hyperparameters $T_{\text{warm}}$, $T_{\text{relabel}}$, buffer size $|\mathcal{B}_{\text{MPC}}|$, mix ratio $p$, label generator $\mathcal{M} \in \{\text{Grad-MPC},\, \text{FMQ}\}$.
\State $\mathcal{B}_{\text{MPC}} \leftarrow \emptyset$.
\For{$\text{step} = 1, \dots, T_{\text{total}}$}
    \If{$\text{step} \ge T_{\text{warm}}$ \textbf{and} $\text{step} \bmod T_{\text{relabel}} = 0$}
        \State Sample $|\mathcal{B}_{\text{MPC}}|$ states from $\mathcal{D}_{\text{off}}$.
        \State For each $z$, set $\mathbf{a}^{\text{MPC}}(z) \leftarrow \mathcal{M}(z;\,\pi_\theta,\,Q_\phi,\,f_\psi)$ using Grad-MPC (\S\ref{sec:mpc-grad}) or FMQ (Alg.~\ref{alg:mpc-fmq}).
        \State $\mathcal{B}_{\text{MPC}} \leftarrow \{(z,\,\mathbf{a}^{\text{MPC}}(z))\}$; reset the actor-only optimiser state.
    \EndIf
    \State Sample minibatch $B$ from $\mathcal{D}_{\text{off}}$.
    \If{$\mathcal{B}_{\text{MPC}} \ne \emptyset$}
        \State Sample $\lfloor pB \rfloor$ pairs from $\mathcal{B}_{\text{MPC}}$ and overwrite the first $\lfloor pB \rfloor$ rows of $B$'s observations and actions.
    \EndIf
    \State Update $(\pi_\theta, Q_\phi)$ on $B$ with the standard FQL critic and actor losses.
\EndFor
\end{algorithmic}
\end{algorithm}

% ---------------------------------------------------------------------
\section{Actor Fine-Tuning with World-Model Transitions}
\label{sec:mpcd-online-wm}

All methods in this chapter so far use the world model either as a
planning tool (Chapter~\ref{chap:wm-planner}) or as a way to relabel
actions in an otherwise real-data training loop
(\S\S\ref{sec:mpcd-offline}--\ref{sec:mpcd-interleaved}). In both cases,
the replay buffer remains grounded in real environment transitions.
This section goes a step further: rather than relabelling actions, we
use the world model to \emph{generate the transitions themselves}, so
that the agent effectively performs offline-to-online reinforcement
learning entirely within the learned dynamics.
This is close to the setting of Chapter~\ref{chap:wm-env}, but with two
additional constraints: actions are selected by a critic-scored MPC
controller, and successor states are generated by the same
action-conditioned model used in the rollout. The design therefore tests
whether synthetic transitions are more useful when action selection and
successor generation are kept consistent.

\paragraph{From action relabelling to full transition synthesis.}
The interleaved method (\S\ref{sec:mpcd-interleaved}) still relies on
the latents from real transitions: it modifies the action in a stored tuple but keeps the
original successor $z'$. This creates a subtle inconsistency: the critic
is trained on tuples of the form
\[
(z,\, \mathbf{a}^{\text{MPC}},\, r,\, z'),
\]
where $z'$ corresponds to the behaviour action $\mathbf{a}$, not the
MPC action $\mathbf{a}^{\text{MPC}}$. The approach works because MPC
keeps $\mathbf{a}^{\text{MPC}}$ close to $\mathbf{a}$, but it does not
represent a true transition under the modified action.

We remove this inconsistency by constructing transitions entirely inside
the world model. Starting from a real initial latent, we use FMQ-MPC to
select an action chunk $\mathbf{a}^{\text{MPC}}$ and then execute it in
the world model:
\[
z'_\psi = f_\psi(z, \mathbf{a}^{\text{MPC}}).
\]
We assign a sparse task reward in latent space,
\[
r_\psi = \mathbb{1}\{\lVert z'_\psi - z_{\text{goal}}\rVert < \delta\},
\]
and obtain a fully self-consistent synthetic transition
\[
(z,\, \mathbf{a}^{\text{MPC}},\, r_\psi,\, z'_\psi).
\]
These transitions are added to a replay buffer seeded with offline data,
and training proceeds as in standard FQL.
The key change is that the successor state is no longer inherited from
the dataset: it is now generated by the same dynamics model that
produced the action-conditioned rollout.

\paragraph{Constraints.}
At first sight, this procedure appears to reintroduce exactly the risks
identified in Chapter~\ref{chap:wm-env}. It is therefore constrained in
two ways.
First, actions are not arbitrary model rollouts but are selected by an
offline-trained critic through MPC. This means the model is only queried
along trajectories that the critic already assigns high value under real
data, effectively restricting synthetic transitions to regions close to
the data manifold.
Second, the world model is only used over short horizons inside MPC
($H=2$), where empirical drift remains limited (\S\ref{sec:wm-drift}).
Thus, the model is never required to produce long-horizon consistency,
only locally plausible one-step transitions under curated actions.

\paragraph{Training regimes.}
We evaluate two variants that differ only in whether the critic is
updated on synthetic transitions.
In the \emph{full update} variant, the critic performs standard FQL
updates on the mixed buffer. This tests whether imagined successors can
safely enter Bellman targets.
In the \emph{frozen critic} variant, the offline-trained critic is kept
fixed throughout the online phase and only used for (i) MPC scoring and
(ii) supervision of the BC-flow actor. This removes value
bootstrapping on imagined transitions. The full comparison is reported
in \S\ref{sec:exp-mpc-online-wm}
(Table~\ref{tab:mpc-online-e2}).

\paragraph{Effect of reintroducing planning at test time.}
We also evaluate whether applying MPC on top of the trained actor still
helps. This checks whether the actor has absorbed the planner's local
improvement signal or whether test-time search remains necessary.

% ---------------------------------------------------------------------
\section{Goal-Conditioned Actor Fine-Tuning with World-Model Transitions}
\label{sec:mpcd-gc}

The single-task setting isolates the online-in-WM loop, but it is also
favourable because the agent only needs to cover one narrow goal
distribution. In the goal-conditioned setting, the policy must represent
a much larger state-goal space, so world-model-driven data collection is
more likely to reduce coverage.
This section evaluates whether the same frozen-critic online-in-WM
procedure from \S\ref{sec:mpcd-online-wm} transfers to the
five-task goal-conditioned setting.

\paragraph{Setup.}
The underlying algorithm is unchanged from
\S\ref{sec:mpcd-online-wm}: FMQ-MPC is used to select actions inside
the world model, the replay buffer is augmented with imagined
transitions, and the critic is frozen so that only the BC-flow actor
is updated.
Three modifications define the goal-conditioned version. First, the
observation is augmented to include the goal embedding:
$[z, z_{\text{goal}}]$. Second, training is performed jointly across a
set of related goal-conditioned tasks rather than a single task. Third,
rewards follow the HIQL HER-style indicator formulation used in the
offline goal-conditioned baseline.
The pure offline goal-conditioned FQL checkpoint serves as the main
reference point for evaluating whether online-in-WM training provides
additional benefit.

\paragraph{HER inside the online-in-WM loop.}
During the online phase, goals can be sampled in two ways.
The default setting uses evaluation goals directly (``eval-goals''),
meaning that the same goal distribution appears both during training
and evaluation. While effective, this introduces a mild evaluation
leakage, since the agent is repeatedly trained on the exact goals it
will later be tested on.
To address this, we also evaluate a hindsight relabelling (HER)
variant \cite{her}, where goals are sampled from achieved states in
previous trajectories. Concretely, each rollout can be relabelled by
treating a future achieved state as the goal, and training goals are
drawn from a buffer of achieved states rather than the evaluation goal
set. In the indicator-reward setting, HER only modifies the goal
argument; the underlying transitions remain unchanged.
Algorithm~\ref{alg:owm-gc} summarises the full procedure.

\begin{algorithm}[h]
\caption{Goal-conditioned online-in-WM with HER. The single-task
controller of \S\ref{sec:mpcd-online-wm} is the special case in
which the goal is fixed across the entire online phase and the HER
relabelling step is skipped.}
\label{alg:owm-gc}
\begin{algorithmic}[1]
\Require offline GC checkpoint $(\pi_\theta, Q_\phi)$, frozen WM $f_\psi$, replay buffer $\mathcal{D}_{\text{off}}$ seeded with the offline transitions, achieved-state buffer $\mathcal{G}$, goal source $\mathcal{S} \in \{\text{eval}, \text{random achieved}\}$.
\State Freeze the critic $Q_\phi$.
\For{$\text{step} = 1, \dots, T_{\text{online}}$}
    \State Sample initial latent $z_0 \sim \mathcal{D}_{\text{off}}$ and goal $g \sim \mathcal{S}$ (or $g \sim \mathcal{G}$ for HER).
    \State Roll out: at each WM step $h$, $\mathbf{a}_h \leftarrow \text{FMQ-MPC}(z_h \oplus g)$ and $z_{h+1} \leftarrow f_\psi(z_h, \mathbf{a}_h)$.
    \State Compute reward $r_h = r_h^{\text{HIQL}}(z_{h+1}, g)$.
    \State Append $(z_h \oplus g,\, \mathbf{a}_h,\, r_h,\, z_{h+1} \oplus g)$ to $\mathcal{D}_{\text{off}}$ and store $z_{h+1}$ in $\mathcal{G}$.
    \State Sample minibatch from $\mathcal{D}_{\text{off}}$, take an actor-only gradient step with $Q_\phi$ frozen.
\EndFor
\end{algorithmic}
\end{algorithm}

The main limiting factor appears to be a combination of reduced
coverage and distributional drift. First, the FMQ-MPC controller
induces a narrow visitation distribution inside the world model, which
progressively biases the online buffer away from the broader
state-action coverage present in the offline data. Second, the
goal-conditioned observation space effectively increases dimensionality
relative to the single-task case, meaning that a fixed number of online
steps explores a smaller fraction of the relevant space.

A useful consistency check arises when comparing with the single-task
setting. There, critic updates on imagined transitions are a risky
design choice because the offline and generated data use different
reward scales. In the goal-conditioned setting, both offline and online
data share the same HIQL-style indicator reward convention, so the same
ablation tests whether critic updates remain fragile once the reward
scale is aligned. The corresponding results are reported in
\S\ref{sec:exp-mpc-online-wm-gc} (Table~\ref{tab:mpc-online-e3}).

\section{Static Dataset Generation with World Model}
\label{sec:mpcd-dataset-gen}

The online-in-WM method of \S\ref{sec:mpcd-online-wm} uses the world
model continuously during training: at every step it generates one MPC
transition, appends it to the replay buffer, and immediately updates
the policy. A natural alternative is to remove this coupling entirely
and use the world model only once, as a static generator of a dataset
that is then used for standard offline training.
This section evaluates whether such a decoupled use of the world model
can reproduce the benefits of online-in-WM training without requiring
any further interaction with the model during learning.

\paragraph{Experimental design.}
We construct two synthetic datasets using the frozen GC checkpoint from
\S\ref{sec:mpcd-gc}, each consisting of full trajectories:

\begin{itemize}
    \item \emph{FMQ rollout}: trajectories are generated by running
    FMQ-MPC inside the world model. At each step, FMQ selects the
    action, and the world model produces the successor state. The goal
    is used only for action selection and is not stored in the dataset.

    \item \emph{BoN rollout}: identical, but without FMQ refinement.
    Actions are taken directly from the best-of-$N$ candidate set
    produced by the BC-flow actor.
\end{itemize}

Both datasets are matched in size and episode structure to the original
offline play dataset (maximum horizon 40 or goal completion). As a
control, we also evaluate the original play dataset under the same HER
goal relabelling procedure, ensuring that all datasets are processed
with identical goal construction at training time.

\paragraph{Evaluation question.}
This experiment tests whether world-model-generated data remains useful
when data generation is decoupled from the learner. In online-in-WM
training, the current actor is repeatedly updated on the same type of
planner-improved actions it helps generate. In the static setting, a
fresh learner receives a fixed corpus generated by another policy and
critic. The comparison is reported in \S\ref{sec:exp-wm-dataset-gen}
(Table~\ref{tab:wm-dataset-gen}).
% ---------------------------------------------------------------------
\section{Summary}
\label{sec:mpcd-discussion}

This chapter defined three ways of turning planner improvements into
training signal: fixed MPC labels, periodically refreshed MPC labels,
and full online fine-tuning inside the world model. The key distinction
is how tightly generated data remains coupled to the current actor and
critic. The detailed results are reported in
\S\S\ref{sec:exp-mpc-interleaved-grad}--\ref{sec:exp-wm-dataset-gen} and
summarised in \S\ref{sec:exp-headline}.
